\chapter{How the Rest of the World Chooses}
\label{c:world}

\begin{glance}
Every chapter so far described a contest between two principals; the outcome will be settled by everybody else. Solar was decided in German feed-in tariffs and desert tenders, batteries in the procurement rules of a dozen auto markets, EVs in Southeast Asia and Latin America: the producer with the best cost curve won because third markets bought it. The non-aligned are the electorate, and their question is not which stack to join but whether they can evaluate, fork, certify and operate components from both while retaining sovereign capability. On the decision matrix the Chinese commons wins on cost, operational sovereignty and export-control exposure; the American stack on frontier capability, trust and enterprise interoperability; and on industrial spillover neither helps unless the adopting country does the work itself. The financed sovereign ``stack'' each principal is said to be exporting has, at the currency date, been delivered by neither anywhere: one contract, self-financed and split, is the whole record.
\end{glance}

\section{The decision matrix}

Table~\ref{tab:decision} lists eight dimensions of a national stack decision, two of them routinely omitted. Export-control exposure is asymmetric but not zero for either stack: American compute and cloud access have been restricted by rule and licence, and Chinese open weights carry a continuity risk---the playbook's endgame is gatekeeping, executed already in solar, batteries and rare earths. Industrial spillover separates a strategy from a purchase: importing tokens builds nothing, while operating models, running probes and contributing to standards builds people, the only input that compounds.

\begin{table}[htbp]
\centering\footnotesize
\caption{Eight dimensions of a national stack decision, and where each stack currently scores better.}
\label{tab:decision}
\begin{tabular}{@{}p{30mm}p{67mm}p{55mm}@{}}
\toprule
\textbf{Dimension} & \textbf{The question} & \textbf{Mid-2026 reading} \\
\midrule
Capability & Is the best available model good enough for the workload? & US stack leads at the absolute frontier; best Chinese open model eight index points behind (v4.3) \\
Cost & Cost per solved task, not per token & Chinese commons (per token, by an order of magnitude; per solved task, unpublished) \\
Sovereignty & Can it be run without egress and without permission? & Chinese commons (self-hostable weights) \\
Trust & Provenance, behaviour, liability, certification & US stack; the commons ships at Level 0--1 \\
Interoperability & Enterprise tooling, integration, harnesses & US stack today; contested at the harness layer \\
Export control & Exposure to restriction by rule, licence, or release policy & Both exposed; commons less so operationally \\
Standards & Who writes the specifications you will depend on? & Both create dependence; China shipping stage-4 defaults \\
Industrial spillover & Does adoption build domestic capability? & Neither---unless the country does the work \\
\bottomrule
\end{tabular}
\end{table}

Applied, the matrix predicts behaviour better than allegiance does: why governments ban a Chinese chat application while their agencies evaluate its weights, why enterprises buy American cloud and Chinese weights simultaneously, and why Chinese models were adopted fastest where the ability to pay rent is least. It also exposes a Chinese vulnerability: the commons wins on cost and operational sovereignty but loses on trust, and trust is what a coalition of non-aligned states could supply for itself---which would make the commons safe to adopt and simultaneously deprive its sponsor of the dependence it purchased.

\section{The export offer, as it exists}
\label{c:export-offer}

The story told about the non-aligned buyer in 2026 is that the choice is no longer between artifacts but between \emph{stacks}---a financed, sovereign-branded Chinese package of chips, datacentres, models and training against an American one, with the telecom campaign of Chapter~\ref{c:precedents} as the template Beijing is supposed to be replaying. Tested against what had been signed at \datafreeze, it is a story. The wrapper is real: the World AI Cooperation Organization was founded in Shanghai on 16~July 2026 by twenty-nine signatories, but its charter contains no compute commitment, no financing clause and no datacentre programme, its dues run on the Universal Postal Union scale, and India, Turkey, Egypt and every Gulf state but Oman stayed away~\cite{waico} [P/A]. The live cases are in Table~\ref{tab:export-offer-c}: one contract, one bid, one evaluation, two frameworks, and no financing from the vendor's side in any row. The American offer fares no better---an executive order, a financing architecture on paper, seventy-eight consortium proposals and zero designations~\cite{us-ai-exports} [P]. \emph{Neither offer has been delivered anywhere}: the contest in Cairo is a bid against a counter-offer for a buyer who has signed with neither, and the buyer that has signed, Brazil, signed with both, financed itself, and said it was doing so in order not to depend on either~\cite{brazil-ai-supercomputer}.

\begin{table}[htbp]
\centering\footnotesize
\caption{The Chinese AI export offer at \datafreeze: no vendor financing in any row~\cite{huawei-egypt,huawei-malaysia,brazil-ai-supercomputer,kazakh-huawei}.}
\label{tab:export-offer-c}
\begin{tabular}{@{}p{16mm}p{52mm}p{38mm}p{38mm}@{}}
\toprule
\textbf{Country} & \textbf{Offer} & \textbf{Status; financing} & \textbf{US response} \\
\midrule
Egypt & Huawei bid: $\sim$2,000 Ascend 950-series ($\sim$1,408 training, $\sim$600 inference), 12-month build & Unsigned; no price; no financing reported; no AI agreement from Xi's Cairo visit & State Dept assembling NVIDIA/AMD/Microsoft counter-offer; warning that Ascend use may breach US rules \\
\addlinespace
Malaysia & Ascend 910C as backbone of a RM2B ($\sim$US\$494M) sovereign project & Under evaluation; unsigned; buyer's own budget & 2025 warning on 910C; export-control terms in the 2025 US trade pact \\
\addlinespace
Brazil & R\$1.3B ($\sim$US\$251M) to Huawei/iFlytek: Rio machine on Ascend plus Portuguese-language model & Contracted 21 Aug.\ 2026; Brazil's FNDCT; second R\$1B machine expected to go to NVIDIA & None reported \\
\addlinespace
Kazakhstan & Huawei partnership; equipment purchase; 1\,GW ``Data Center Valley'' & MoUs; values undisclosed & None reported \\
\addlinespace
Pakistan (Punjab) & Six datacentres; 1,000 ``AI master trainers'' & Agreement 28 Aug.; no chips or money described & None reported \\
\bottomrule
\end{tabular}
\end{table}

The telecom template fails on its three load-bearing elements. \emph{No credit line}: no CDB or Export--Import Bank facility for AI exports has surfaced---an absence in the record, not a verified absence---and the one contract was financed by the recipient. \emph{No surplus capacity}: Huawei exported switches from factories the state wanted running, whereas DeepSeek's reported order for at least 160,000 Ascend 950DT eats up to a year of output, a reported requirement that state-funded compute run at least 80\% domestic silicon pulls every chip home before any can be spared for Cairo, and BIS guidance holds that Ascend use anywhere risks violating US controls~\cite{deepseek-ascend-order,ascend-supply-ceiling} [A/C]. \emph{And the artifact locks in nothing}: the diffusion mechanism that actually operates is a download---151,448 Qwen derivatives on Hugging Face, over five times Llama's monthly volume~\cite{qwen-derivatives} [V]---which costs the recipient nothing and binds it to nothing; the buyer that will not sign for Ascend runs Qwen already. Three things cut against the conclusion and are printed with it: intent precedes capacity in every precedent, the CDB facility of 2004 having been opened for a company that had been reselling imported switches seventeen years earlier; Egypt is the first Ascend export \emph{attempt} on the record and the live test this section will be scored against; and Brazil is a real contract, on Chinese silicon, in a G20 country's language. One contract is not a campaign; it is the first row of the table a campaign would be measured in \full{22} \full{6}.

\section{Nine positions}

\textbf{Japan} holds the strongest technical hand of any non-principal and the weakest demand-side position: exascale HPC lineage, semiconductor materials and equipment franchises, leading-edge memory and packaging, the only measurement-anchored national demand model for AI for science, and standing in both blocs; against that, a model ecosystem positioned for Japanese-language enterprise depth rather than the open frontier, an energy position that cannot win a watts race, and research provisioning an order of magnitude below its own modeled requirement. Its 2025--26 record~\cite{japan-ai-act} carries one honest tension: FugakuNEXT's architecture buys NVIDIA coupling at the exact moment the compute chapter argues the accelerated-CPU corner is the sovereign long game---defensible as a hedge, dangerous as a habit. \textbf{The European Union} has the pen and not the press, and in 2025--26 the pen hesitated: high-risk AI Act obligations pushed back by up to two years~\cite{eu-omnibus} while thirteen AI factories and an oversubscribed tender for up to seven ``gigafactories'' went ahead~\cite{eu-gigafactories,eu-giga-tender}---seven empty cathedrals of compute would be the solar-tariff mistake with better architecture. \textbf{India} is the largest undecided market, absent from WAICO, and its national mission's roughly 38,000 subsidized GPUs at under a dollar an hour with utilization lagging badly~\cite{indiaai} prove that compute is the easy term and integration the binding one. \textbf{South Korea} is both a critical supplier---SK hynix holds about 62\% of HBM~\cite{korea-hbm}---and a sovereign builder whose state-funded foundation-model tournament with open-weight release as the qualifying bar is playbook step four executed by a democracy. \textbf{Singapore and ASEAN} face the small-state problem, solved by autonomy through selection rather than scale; ASEAN is the clearest third-market dynamic, adopting Chinese models fastest for the reasons it adopted Chinese EVs, and Malaysia is the live case of what a US trade pact does to an Ascend evaluation. \textbf{The Gulf} is capital-rich, energy-rich and compute-poor by choice: the swing buyer in any repricing scenario, the first place Washington has located frontier-scale capacity it controls only contractually~\cite{gulf-ai}, and the market where Huawei's 2025 pitch of low thousands of 910B closed nothing. \textbf{Taiwan} is the busbar chapter compressed into one island. \textbf{Brazil and Latin America} face near-pure commons adoption plus regional pooling; Brazil's split, self-financed programme is the one executed template in the record. \textbf{Africa} is a recipient in most Western strategy documents and a partner in Chinese AI diplomacy: the fastest-growing developer base in the world, no real access track built for it, WAICO bidding, and Egypt the test of whether any Ascend export can close \full{22}.

Chapter~\ref{c:mirror} changes the menu the matrix was drawn for. It was drawn for a Chinese commons, free to download but unaudited and jurisdictionally awkward, against a Western product, metered but certified and integrated. The buyer now chooses between \emph{two commons}, and the second arrives with its training data attached and a chip vendor's balance sheet behind it~\cite{nemotron,amd-open} [P]. But the second has a single principal contributor with a direct commercial interest in the complement, so \emph{openness supplied by the vendor of the complement is openness with a direction}---free at the model layer precisely because it is not free at the silicon layer, as the sponsor says on the record. That is the mechanism of \S\ref{c:theory} working as predicted, and the practical consequence is that the durability of this openness is a function of the sponsor's strategy rather than of the licence, and the place to monitor it is an earnings call. Asking about the licence is the least informative question available; ask instead about the \emph{provenance of the openness}---who funded the training run and at which layer they earn; whether the training data is released, so the artifact can be re-derived rather than merely fine-tuned; what happens to the \emph{next} release if the sponsor's strategy changes; whether the distribution channel is independent of both sponsor and complement vendor, and whether you hold a mirror that survives its loss; and whether anyone who does not earn on the answer has evaluated the model. The last is the question no vendor in either commons can answer, which is why two competing commons raise the supply of free artifacts and do nothing for the supply of independent assurance about them.

\section{Five moves for any non-aligned country}

Regardless of size: measure your own demand in a recognized unit (Chapter~\ref{c:science}); adopt both commons and audit both, with the full trust stack; retain a pretraining capability sized to competence rather than to a benchmark; pool internationally for capacity and certification; and do the integration work domestically, because that is the only place spillover accrues. The trust gap is closable with infrastructure rather than prohibition, and the window is quarters rather than decades: the standards layer is being set now and the certification layer is still vacant.

